\section{ANÁLISE DOS RESULTADOS}

\emph{A. Fotodiodo}

Com os valores de tensâo e corrente da tabela (I) foi obtida a curva carcterísitca (7). Pelo gráfico fica perceptível que o potencial de contato do diodo está em torno dos 1,9 V, já que é apartir desse valor que começa o fluxo de corrente. Com o aumento da tensâo da fonte percebe-se um aumento da corrente ID, e, consequentemente, o diodo brilhava mais assim como esperado.

%  gráfico - Curva I_D x V_D

\emph{B. Fototransistor}

A partir dos dados obtidos experimentalmente, vide tabelas
II e III, obtivemos as curvas apresentadas na figura abaixo:

%  gráfico - Curva I_C x V_CE

Analisando a figura 8 podemos notar que o aumento da luminosidade, através da tensâo do varivolt, influencia diretamente na variaçâo (aumento) da corrente de saturaçâo IC, visto que a corrente de saturaçâo, com o varivolt em 60 mV, é aproximadamente 0,12 mA e, com o varivolt em 80 mV, a corrente de saturaçâo é cerca de 0,62 mA.

Assim, podemos dizer que a corrente no coletor e no emissor é controlada pela intesidade luminosa incidida sobre o fototransistor.

\emph{C. Fotocélulas}

A tensâo de circuito aberto e a corrente de curto-circuito
sâo medidas importantes que caracterizam a fotocélula, ambas
grandezas aumentam com a intensidade luminosa incidida
sobre a célula, assim como mostra a tabela IV.

%  gráfico - Curva I_sc x Voc

A curva apresentada na figura 9 mostra que a corrente de
curto circuito aumenta exponecialmente em relaçâo a tensâo
de circuito aberto.

\emph{D. Acoplamento óptico}

Q1: Porque o sinal da saida está defasado ao respeito do sinal da entrada?

A defasagem (de 90º) do sinal de saída ocorre devido ao tempo em que o sinal de entrada é produzido e passa pela LED e é captado pelo fototransistor.

Q2: Porque quando a amplitude da onda de entrada é aumentada a amplitude da onda de saida é aumentada e a amplitude da onda saída é cortado?

A amplitude do sinal de saída cresce junto com a resistencia do fototransistor que cresce com o aumento da amplitude do sinal de entrada. Mas esse crescimento é limitado pela tensâo V2, quando a amplitude começa a ultrapassar V2 ela é cortada e passa a ter um formato reto nos picos da onda.

Q3: Porque em altas frequencias há uma diminuiçâo na amplitude da onda de saída e também, uma defasgem?

Essa diminuiçâo é devida a limitaçâo da velocidade de resposta do fototransistor. Tanto o LED quanto o fototransistor possuem um limite de frequência.


\vspace{.25cm}
